%%%%%%%%%%%%
%
% $Autor: Wings $
% $Datum: 2019-03-05 08:03:15Z $
% $Pfad: SafetyGuidelines.tex $
% $Version: 4250 $
% !TeX spellcheck = en_GB/de_DE
% !TeX encoding = utf8
% !TeX root = manual 
% !TeX TXS-program:bibliography = txs:///biber
%
%%%%%%%%%%%%

\chapter{Safety Guidelines}

\section{Development Environment Safety}

\begin{itemize}
    \item Never modify core localization algorithms without testing in simulator first
    \item Always implement emergency stop handlers in new code
    \item Use feature flags (configuration options) rather than code comments for experimental features
    \item Include error handling and graceful degradation for sensor failures
\end{itemize}

\section{Testing Safety Practices}

\begin{itemize}
    \item Create unit tests that verify safety guarantees (e.g., max acceleration limits)
    \item Test error recovery paths, not just happy paths
    \item Simulate sensor failures and verify appropriate robot response
    \item Use safety limits (max speed, max turn rate) in all motion commands
\end{itemize}

\section{Code Review Standards}

\begin{itemize}
    \item All changes to motion control require at least two approvals
    \item Safety-critical sections must have comprehensive comments
    \item Require evidence of both unit testing and hardware testing
    \item Document any relaxing of safety constraints and provide justification
\end{itemize}

\section{Performance and Reliability}

\begin{itemize}
    \item Monitor real-time thread priorities; ensure safety loops run at highest priority
    \item Implement watchdog timers to detect system hangs
    \item Test code under high CPU load to verify safety features remain responsive
    \item Use mutex and spinlock carefully to avoid deadlocks that could leave motors running
\end{itemize}





